\documentclass[11pt,a4paper]{article}

% Works with pdfLaTeX and Overleaf.
\usepackage[a4paper,margin=24mm]{geometry}
\usepackage[T1]{fontenc}
\usepackage[utf8]{inputenc}
\usepackage{lmodern}
\usepackage{microtype}
\usepackage{graphicx}
\usepackage{booktabs}
\usepackage{array}
\usepackage{amsmath}
\usepackage{enumitem}
\usepackage{hyperref}
\usepackage{xurl} % allow line breaks in long dataset URLs
\usepackage{float}
\usepackage{caption}
\usepackage{xcolor}
\hypersetup{colorlinks=true,linkcolor=blue!55!black,urlcolor=blue!55!black,citecolor=blue!55!black}
\setlength{\parindent}{0pt}
\setlength{\parskip}{6pt}
\renewcommand{\arraystretch}{1.16}
\captionsetup{font=small}

\title{\textbf{Probability Iterative Closest Point: 2D Reproduction and 3D Registration}\\[4pt]
       \large Computer Vision -- Research Paper Reproduction Project}
\author{\textbf{Team members:}\\
        Student 1 -- Full name\\
        Student 2 -- Full name\\
        Student 3 -- Full name\\[5pt]
        Vrije Universiteit Brussel}
\date{Submission deadline: 21 December 2026}

\begin{document}
\maketitle
\vspace{-4mm}

\begin{center}
\textbf{Reproduced article:} J. Liu, S. Du, C. Zhang, J. Zhu, K. Li, and J. Xue,
``Probability Iterative Closest Point Algorithm for Position Estimation,''
\emph{2014 IEEE 17th International Conference on Intelligent Transportation Systems},
pp. 458--463, 2014. DOI: \href{https://doi.org/10.1109/ITSC.2014.6957732}{10.1109/ITSC.2014.6957732}.
\end{center}

\section{Summary}
We studied Probability Iterative Closest Point (PICP), a way to align point
sets when some points are noisy or misplaced. We reproduced the main idea on
three 2D shapes, then extended the comparison to three-dimensional data. The
3D datasets include two generated surfaces and the \textbf{Stanford Bunny},
a reconstructed model from real optical scans. We compared ICP, a simplified
Soft-EM-ICP method, PICP with the original project settings, and PICP with
slower annealing. To evaluate the methods fairly, we used the same test cases
and measured nearest-neighbour error, known inlier error, and 3D rotation and
translation errors. PICP often aligned the known inlier points more accurately
under synthetic outlier corruption. It did not always achieve the smallest
all-point error and could struggle with some starting positions or incomplete
shapes. On the Bunny, we applied known artificial transformations and noise
to the real reconstructed geometry. Therefore, these results show controlled
registration on real 3D geometry, \textbf{not} alignment of two independently
captured scans. The new work extends the paper and does not directly reproduce
its numerical results.

\section{Implementation}

\subsection{Data acquisition}
We used the \textbf{MPEG-7 CE-Shape-1 Part B} shape dataset, which is also used
in the original PICP paper~\cite{liu2014}. The full dataset contains 1,400
binary images in 70 shape groups~\cite{dataset}. For this project we used three
included images: \texttt{apple-1.gif}, \texttt{cup-1.gif}, and
\texttt{children-1.gif}. They are provided in the project's \texttt{data/}
folder, so no download is needed for the main experiment.

Dataset information and a download source are available on the
\href{https://www.ehu.eus/ccwintco/index.php/MPEG-7_Core_Experiment_CE-Shape-1_Test_Set._Benchmarking_image_database_for_shape_recognition_techniques}{MPEG-7 dataset information page}.
The original paper is included in our project as \texttt{PICP.pdf}.

\subsection{Preprocessing}
Each GIF image is converted into a set of 2D contour points:
\begin{enumerate}[leftmargin=*,itemsep=2pt]
  \item Open the image in grayscale and use intensity 128 to separate the shape from its background.
  \item Find boundary pixels using horizontal and vertical changes in the binary image.
  \item Convert image coordinates to $(x,-y)$ and move the shape centre to $(0,0)$.
  \item Keep at most 600 boundary points, selected at evenly spaced indices from the detected points.
\end{enumerate}
We do not resize the images or change their scale. We then make a second,
misaligned point set by rotating the shape by $15^\circ$ and adding a translation
of $(12,-8)$ pixels. We add Gaussian noise with standard deviation 0.4 pixels.
We also add a stronger displacement with mean $(-12,-12)$ pixels and standard
deviation 2 pixels to points in a selected spatial region. These points are
marked as synthetic outliers. The known point correspondence is kept so that
we can measure inlier error accurately.

The main experiment uses \textbf{five trials per shape} and \textbf{600 points}
(where available). Its base random seed is \textbf{42}; individual seeds are
computed as $42 + 1000s + t$, where $s$ is the shape index and $t$ is the
trial number starting from zero. Thus, repeated runs create the same cases.

\subsection{New 3D data and preprocessing}
For the 3D extension, we first created two non-planar point clouds:
\texttt{asymmetric\_shell.ply} and \texttt{curved\_tube.ply}. These are
\textbf{synthetic} shapes. The original code and seed 2026 are provided so they
can be generated again. We also added the high-resolution \textbf{Stanford
Bunny}, \texttt{bun\_zipper.ply}, which has 35,947 vertices. This mesh was
reconstructed from real optical scans at Stanford University Computer Graphics
Laboratory~\cite{stanfordrepo}. It is included in \texttt{data/3d/}, together
with one original raw scan for provenance. Our experiments use the
\emph{reconstructed mesh vertices}, not the raw scan. We acknowledge Stanford
as the source of this dataset.

For each 3D object, we centre the point coordinates and divide by the longest
bounding-box side, so the longest dimension is one unit. We select 350 points
with a fixed seed. We then create an artificial source by rotating these
points by Euler xyz angles $(12^\circ,-9^\circ,15^\circ)$ and translating by
$(0.18,-0.10,0.08)$. We add Gaussian noise with standard deviation 0.005
normalised units and displace 10\% or 30\% of source points as outliers.
Outlier displacement has mean $(0.22,-0.22,0.15)$ and standard deviation
0.18. We use three fixed trials per condition, with base seed 20261112.

In a harder 3D test we keep only the lower 75\% of the source points, based
on the $z$ coordinate, while leaving the target surface complete. Because we
created the misalignment ourselves, we know the original matching point and
the correct inverse rotation and translation. This allows direct error
measurement. Even for the Stanford Bunny, the registration pairs are
\textbf{synthetically transformed copies of the same scanned reconstruction},
not two separate camera or scanner captures.

\subsection{Algorithms}
The task is \emph{rigid point-set registration}: find a rotation $R$ and a
translation $t$ so that transformed source points align with model points.
In 2D, $R$ is a $2\times 2$ matrix; in 3D, it is a $3\times 3$ matrix:
\begin{equation}
  x_{\mathrm{aligned}} = Rx + t.
\end{equation}
All three methods repeat two basic steps: find corresponding points and update
the transformation.

\textbf{ICP:} Each source point is matched to its nearest model point. An SVD
calculation finds the rotation and translation that reduce the squared distance
between the matched points. All points have equal influence.

\textbf{Soft-EM-ICP (simplified):} Each source point has a soft match to several
model points. These matches are used to find a new transformation. We use a
simplified soft-correspondence implementation, not the full EM-ICP method from
the original paper.

\textbf{PICP:} As in ICP, each source point finds its nearest model point.
However, PICP also gives each point a probability weight. Points with large
matching errors can receive smaller weights, so they influence the next SVD
transformation less. PICP updates the Gaussian variance during its iterations.
This follows the main idea in the paper~\cite{liu2014}.

\subsection{Implementation choices and differences from the paper}
The code is written in Python with NumPy, SciPy, Pillow, Matplotlib, and Pandas.
The same weighted SVD routine supports both 2D and 3D coordinates. The 3D
extension uses SciPy's 3D rotation utilities, an ASCII/binary PLY loader,
and Matplotlib's 3D plotting tools. Open3D is optional and is not required
to reproduce the saved figures.
We use a \texttt{KDTree} to find nearest neighbours; the original paper uses
Delaunay triangulation for this search. The basic nearest-point rule is the
same, but the search method is different. The paper describes an EM-ICP
comparison; our Soft-EM-ICP is a simplified baseline. Our noise generator and
outlier percentages are also our own controlled tests. Finally, we did not
reproduce the paper's real vehicle-position experiment, because that dataset
is not part of this project.

For the main experiment, PICP uses annealing parameter $\lambda=1.4$ and initial
variance scale $0.20$. ICP and PICP allow up to 60 iterations. A later
\emph{exploratory} experiment also tests $\lambda=1.1$ and scale $0.05$.
These alternative settings are reported separately rather than silently
replacing the main settings.

\subsection{Environment and runtime}
This project runs on a CPU; a GPU and model training are not required.
The Python packages are listed in \texttt{requirements.txt}. The main experiments run with \texttt{python run\_all.py} from the project folder:
2D main, synthetic 3D full overlap, Stanford Bunny full overlap, both 75\%
source-coverage partial-overlap experiments, and 2D table generation.
The Bash helper \texttt{run.sh} can also be used where Bash is installed.
The robustness and
sensitivity experiments are separate Python scripts described in the README.

\textbf{Complete before submission:} CPU: \underline{\hspace{38mm}};
RAM: \underline{\hspace{20mm}} GB;
GPU: not used;
main experiment runtime: \underline{\hspace{20mm}} seconds.
These values should be measured on the computer used for the submitted runs.

\subsection{Optional single-case speed demonstration}
We also provide \texttt{benchmark\_single\_case.py} to compare the execution time
and registration error of the methods on \emph{one fixed input case}. The default
examples are a 2D Apple contour (300 points) and a 3D Stanford Bunny sample
(350 points). All methods are given the same source and target points for each
case. We time only the registration call, not data loading, evaluation, or
plotting. The script outputs per-method runtime, number of iterations,
all-point nearest-neighbour RMSE, direct inlier RMSE, and PICP's internal
weighted RMS, together with alignment figures. The resulting CSV is saved at
\path{results/single_case_benchmark/single_case_results.csv}.

This is an \textbf{optional demonstration}, not part of the multi-trial tables
below. We do not report its runtime values as fixed reference numbers: runtime
depends on the machine and one timed run is not enough for a reliable speed
ranking. For a strong timing comparison we would repeat the same case several
times and report the median and variation. The 2D and 3D error units are also
different and should not be directly compared.

\section{Results}

\subsection{How we measure alignment}
We report the \textbf{nearest-neighbour root mean square error (NN-RMSE)} for
all methods. This is the average distance from each aligned point to its
nearest model point, expressed in pixels. A lower value means the points are
closer together. We also report \textbf{direct inlier RMSE}, which compares
uncorrupted points to their known original corresponding points. This is
possible because we created the test data ourselves. These two measures answer
different questions and should not be mixed up.

PICP also computes its own \emph{probability-weighted RMS}. Its small value
must \textbf{not} be treated as directly better than ICP's unweighted RMSE.
Tables below show mean $\pm$ standard deviation over five trials.

\subsection{Main quantitative results}
\begin{table}[H]
\centering
\caption{All-point nearest-neighbour RMSE in pixels (lower is better).}
\label{tab:allrmse}
\small
\begin{tabular}{lccc}
\toprule
Shape & ICP & Simplified Soft-EM-ICP & PICP \\
\midrule
Apple & $4.747\pm0.047$ & $4.748\pm0.047$ & $5.857\pm0.028$ \\
Cup & $4.801\pm0.039$ & $4.805\pm0.039$ & $5.704\pm0.036$ \\
Children & $4.831\pm0.046$ & $4.837\pm0.049$ & $5.497\pm0.068$ \\
\bottomrule
\end{tabular}
\end{table}

In Table~\ref{tab:allrmse}, standard ICP has the lowest all-point NN-RMSE
for all three shapes. This metric includes the corrupted outlier points.

\begin{table}[H]
\centering
\caption{Direct inlier RMSE in pixels using known synthetic point correspondences (lower is better).}
\label{tab:directrmse}
\small
\begin{tabular}{lccc}
\toprule
Shape & ICP & Simplified Soft-EM-ICP & PICP \\
\midrule
Apple & $6.876\pm0.138$ & $6.735\pm0.076$ & $4.956\pm0.222$ \\
Cup & $4.776\pm0.062$ & $4.944\pm0.071$ & $0.576\pm0.019$ \\
Children & $3.375\pm0.083$ & $3.272\pm0.076$ & $0.565\pm0.018$ \\
\bottomrule
\end{tabular}
\end{table}

In Table~\ref{tab:directrmse}, PICP gives a lower direct inlier error for
every shape. Its result is especially strong on Cup and Children. However,
this does not mean it is the best method on every measure.

\subsection{Qualitative results}
Figure~\ref{fig:input} shows one synthetic registration case. It contains
a shifted and rotated source shape, with some noisy points.

\begin{figure}[H]
\centering
\includegraphics[width=0.72\textwidth]{figures/figure1_synthetic_input.png}
\caption{Example input point sets from the Apple experiment, including synthetic outliers.}
\label{fig:input}
\end{figure}

Figure~\ref{fig:registration} shows the registered Apple point sets for the
three methods. It is useful to compare the overall shape and the points that
are still far from the model. Figure~\ref{fig:convergence} shows how the
common nearest-neighbour error changes over iterations.

\begin{figure}[H]
\centering
\includegraphics[width=\textwidth]{figures/figure2_apple_registration.png}
\caption{Apple registration results: ICP, simplified Soft-EM-ICP, and PICP.}
\label{fig:registration}
\end{figure}

\begin{figure}[H]
\centering
\includegraphics[width=0.82\textwidth]{figures/figure3_common_metric_convergence.png}
\caption{Convergence on Apple, evaluated with the same all-point NN-RMSE metric.}
\label{fig:convergence}
\end{figure}

\subsection{Additional tests: outliers and PICP settings}
We also tested 10\%, 20\%, 30\%, and 40\% randomly selected outliers. Each
setting used five trials and 300 points. For example, at 30\% outliers, the
following direct inlier errors were observed:

\begin{table}[H]
\centering
\caption{Robustness test at 30\% synthetic outliers: direct inlier RMSE (mean $\pm$ standard deviation, pixels).}
\small
\begin{tabular}{lccc}
\toprule
Shape & ICP & Simplified Soft-EM-ICP & PICP \\
\midrule
Apple & $6.646\pm0.289$ & $5.309\pm0.240$ & $8.478\pm4.047$ \\
Cup & $4.391\pm0.746$ & $6.022\pm0.131$ & $1.444\pm1.946$ \\
Children & $4.727\pm0.256$ & $4.747\pm0.113$ & $0.571\pm0.027$ \\
\bottomrule
\end{tabular}
\end{table}

This test shows that PICP's advantage depends on the shape and the noise.
In particular, its Apple result is unstable at 30\% outliers.

To investigate Apple, we tested 12 combinations of annealing and initial
variance settings. We selected the slower-annealing setting for an
\emph{exploratory} follow-up test on new random seeds; it was not selected
before looking at the Apple results.

\begin{table}[H]
\centering
\caption{Extra 600-point validation: direct inlier RMSE on five new trials (pixels).}
\label{tab:sensitivity}
\small
\begin{tabular}{lccc}
\toprule
Shape & ICP & PICP main setting & PICP exploratory setting \\
\midrule
Apple & $6.978\pm0.317$ & $4.893\pm0.278$ & $1.821\pm0.081$ \\
Cup & $4.727\pm0.095$ & $0.583\pm0.011$ & $1.613\pm0.068$ \\
Children & $3.376\pm0.070$ & $0.589\pm0.015$ & $0.573\pm0.006$ \\
\bottomrule
\end{tabular}
\end{table}

Slower annealing improves Apple substantially, but it makes Cup worse.
Therefore, one set of parameters is not automatically best for every shape.
We also tested different starting angles on Apple. With a $30^\circ$ initial
rotation, the direct inlier error increased to 42.151 pixels for the main PICP
setting and 38.873 pixels for the slower-annealing setting. This is a failure
case: both methods can struggle when the starting alignment is too far away.

\begin{figure}[H]
\centering
\includegraphics[width=0.84\textwidth]{figures/apple_parameter_sensitivity.png}
\caption{Apple parameter sweep. This diagnostic checks how PICP responds to different annealing and variance settings.}
\label{fig:sensitivity}
\end{figure}

\subsection{New 3D results: full overlap}
The two synthetic 3D shapes are different from the MPEG-7 2D shapes. Here,
all distances use \textbf{normalised object units}, not pixels. We measured
the direct inlier RMSE (using known original matching points), the common
all-point NN-RMSE, rotation error in degrees, and translation error.
Rotation error is the angle between the estimated inverse pose and the known
inverse pose. Translation error is their Euclidean distance.
Each table entry below is a mean over three fixed trials.

\begin{table}[H]
\centering
\caption{3D experiment with 30\% injected outliers and full overlap: mean errors (lower is better).}
\label{tab:3d}
\scriptsize
\begin{tabular}{llrrrr}
\toprule
3D dataset & Method & Direct inlier RMSE & All-point NN-RMSE & Rotation ($^\circ$) & Translation \\
\midrule
Shell & ICP & 0.1575 & 0.1424 & 14.63 & 0.1404 \\
Shell & Soft-EM & 0.1486 & 0.1428 & 14.30 & 0.1300 \\
Shell & PICP main & 0.0086 & 0.1623 & 0.13 & 0.0010 \\
Shell & PICP slow & 0.0086 & 0.1619 & 0.15 & 0.0013 \\
Curved tube & ICP & 0.1388 & 0.1433 & 7.96 & 0.1298 \\
Curved tube & Soft-EM & 0.1149 & 0.1445 & 9.39 & 0.1021 \\
Curved tube & PICP main & 0.0089 & 0.1673 & 0.19 & 0.0015 \\
Curved tube & PICP slow & 0.0089 & 0.1668 & 0.12 & 0.0012 \\
\bottomrule
\end{tabular}
\end{table}

In Table~\ref{tab:3d}, PICP has lower \emph{direct inlier error} and lower
rotation error than ICP in these synthetic cases. But PICP's all-point
nearest-neighbour error can be higher: the injected outliers remain far from
the original surface. These two measures answer different questions. The
slower-annealing version was selected after earlier 2D Apple testing, so it
is reported as an exploratory comparison, not as a new paper-prescribed rule.

\begin{figure}[H]
\centering
\includegraphics[width=0.95\textwidth]{figures/3d_asymmetric_shell_30pct_registration.png}
\caption{3D asymmetric-shell example with 30\% synthetic outliers. Blue is the target and red is the source; the panels compare before registration and four methods.}
\label{fig:3d_full}
\end{figure}

\begin{figure}[H]
\centering
\includegraphics[width=0.70\textwidth]{figures/3d_curved_tube_30pct_convergence.png}
\caption{3D curved-tube common NN-RMSE during registration. All algorithms use the same measure on this graph.}
\label{fig:3d_curve}
\end{figure}

\subsection{3D partial-overlap test}
We also removed the upper 25\% of each source surface, while keeping the full
target, and added 30\% outliers. This is a harder synthetic case. In Table
\ref{tab:3d_partial} the original PICP sometimes fails to recover the pose.
The slower-annealing PICP often performs better, but the small number of
trials means this is an exploratory finding.

\begin{table}[H]
\centering
\caption{3D partial overlap (75\% source coverage, 30\% outliers): mean direct inlier error and pose success.}
\label{tab:3d_partial}
\small
\begin{tabular}{llrr}
\toprule
Dataset & Method & Direct inlier RMSE & Pose success \\
\midrule
Shell & ICP & 0.1299 & 0\% \\
Shell & Soft-EM & 0.2096 & 0\% \\
Shell & PICP main & 0.0261 & 67\% \\
Shell & PICP slow & 0.0087 & 100\% \\
Curved tube & ICP & 0.1855 & 0\% \\
Curved tube & Soft-EM & 0.1824 & 0\% \\
Curved tube & PICP main & 0.0976 & 0\% \\
Curved tube & PICP slow & 0.0183 & 67\% \\
\bottomrule
\end{tabular}
\end{table}

\begin{figure}[H]
\centering
\includegraphics[width=0.92\textwidth]{figures/3d_partial_curved_tube_30pct_registration.png}
\caption{Partially visible curved tube with 30\% synthetic outliers. Missing source points make registration more difficult.}
\label{fig:3d_partial}
\end{figure}

A successful pose has rotation error below $5^\circ$ \emph{and} translation
error below 0.05 normalised units. This is our evaluation threshold, not a
threshold from the original PICP paper. The CSV files give results for
both 10\% and 30\% outliers, including standard deviations, runtimes, and
effective sample sizes. They are stored in \texttt{results/3d/} and
\texttt{results/3d\_partial/}.

\subsection{Stanford Bunny: scanned 3D geometry}
We also tested the four methods on vertices sampled from the
\textbf{real Stanford Bunny reconstructed model}. The data were scanned
optically and reconstructed by Stanford~\cite{stanfordrepo}. The test itself
still uses a \emph{synthetically generated source}: the model vertices are
transformed, then affected by simulated noise and outliers. This controlled
test should not be described as independent scan-to-scan registration.

Table~\ref{tab:bunny} reports the mean over three fixed trials with 350
sampled target points, full overlap, and 30\% artificially corrupted source
points. Errors use normalised object units, not pixels. PICP is more accurate
on the known clean correspondences and the true 3D pose, while ICP has a
smaller all-point NN-RMSE because this metric includes outliers.

\begin{table}[H]
\centering
\caption{Stanford Bunny reconstruction, 30\% simulated outliers and full overlap: average of three trials (lower is better).}
\label{tab:bunny}
\small
\begin{tabular}{lrrrr}
\toprule
Method & Direct inlier RMSE & All-point NN-RMSE & Rotation ($^\circ$) & Translation \\
\midrule
ICP & 0.1359 & 0.1355 & 12.79 & 0.1139 \\
Soft-EM (simplified) & 0.1278 & 0.1360 & 8.18 & 0.1185 \\
PICP original & 0.0086 & 0.1491 & 0.24 & 0.0010 \\
PICP slow & 0.0086 & 0.1488 & 0.20 & 0.0013 \\
\bottomrule
\end{tabular}
\end{table}

\begin{figure}[H]
\centering
\includegraphics[width=0.96\textwidth]{figures/3d_bunny_30pct_registration.png}
\caption{Real Stanford Bunny reconstructed geometry under a simulated rigid transform, noise, and 30\% outliers. Blue points are the model; red points show the source before and after alignment.}
\label{fig:bunny}
\end{figure}

We also used 75\% source coverage with 30\% outliers. Table~\ref{tab:bunny_partial}
shows that PICP remained accurate in the three fixed trials. These are
promising controlled tests, but there are too few trials to draw a broad
conclusion for arbitrary real-world scans.

\begin{table}[H]
\centering
\caption{Stanford Bunny with 75\% simulated source coverage and 30\% outliers: three-trial means. Pose success uses the threshold defined earlier.}
\label{tab:bunny_partial}
\small
\begin{tabular}{lrr}
\toprule
Method & Direct inlier RMSE & Pose success \\
\midrule
ICP & 0.0729 & 33\% \\
Soft-EM (simplified) & 0.0813 & 0\% \\
PICP original & 0.0086 & 100\% \\
PICP slow & 0.0087 & 100\% \\
\bottomrule
\end{tabular}
\end{table}

The full numerical results, including standard deviations, run times,
effective sample sizes, and rotation/translation errors, are in
\texttt{results/3d\_bunny/} and \texttt{results/3d\_bunny\_partial/}.

\subsection{Comparison with the original paper}
The original paper reports the following RMS values in its Table II:

\begin{table}[H]
\centering
\caption{RMS values reported by Liu et al.~\cite{liu2014}, Table II. These are \emph{not} our results.}
\small
\begin{tabular}{lccc}
\toprule
Shape & ICP & EM-ICP & PICP \\
\midrule
Apple & 7.3768 & 7.2094 & 0.0035 \\
Cup & 4.3299 & 3.9069 & 0.0047 \\
Children & 1.4703 & 1.4144 & 0.0057 \\
\bottomrule
\end{tabular}
\end{table}

The paper finds that PICP performs very well on its experiments. Our tests
also show an improvement in direct inlier alignment for the main synthetic
cases, but our all-point NN-RMSE does not show the same pattern. We must
\textbf{not compare the numbers in the paper's table directly with ours}:
the corruption process, number of points, implementation details, and reported
RMS measures are not confirmed to be identical.

\section{Discussion}
PICP can help when some points are unreliable because it gives different
weights to different points. In our main experiment this helped PICP align
the unchanged points, especially for Cup and Children. However, outliers still
count in the all-point NN-RMSE, which was higher for PICP. This is why we
show both evaluation measures. The new 3D experiments show a clearer
advantage for PICP on known clean correspondences when many outliers are
present, but these are synthetic objects with synthetic corruption. The
partial-overlap experiment illustrates that some PICP settings can still
fail and that parameter choice matters.

We also noticed that the main PICP setting gave a very small \emph{effective
sample size}: about 5 points on Apple, 15 on Cup, and 4 on Children, even
though the original point sets contain hundreds of points. This means that
a small number of points carry most of the probability weight. Slower
annealing can reduce this problem, but it does not improve every shape.

Our results are different from the paper partly because we used a custom
synthetic noise model, only three example shapes, a simplified EM-ICP
baseline, and a KD-tree implementation. We did not have the vehicle-position
data used in the second part of the paper. Another limitation is that PICP
can fail when the initial rotation is large. A larger evaluation on more
shapes and a closer match to the paper's original settings would improve
this reproduction. The 3D extension includes two generated objects and the real Stanford Bunny
reconstruction, each with three trials per condition. We use synthetic poses
and noise even on the Bunny, so we have not tested alignment of two independent
scanner captures. We cannot claim universal performance on real-world scans.
Rotation and translation error are based on synthetic known poses, not
inferred from actual sensor records.

\textbf{Conclusion:} We implemented the main idea of PICP and compared it
with standard ICP and a simplified Soft-EM-ICP baseline in 2D and 3D. PICP
showed strong inlier alignment on the synthetic surfaces and on synthetic
transformations of Stanford Bunny scan geometry. It did not always have the
lowest all-point nearest-neighbour error and could fail with partial overlap. The new work addresses 3D registration while keeping the
original 2D reproduction available. It is not evidence of universal
superiority in registration of independent real 3D sensor scans.

\section{Reproducibility checklist}
\begin{itemize}[leftmargin=*,itemsep=3pt]
\item \textbf{Code:} Included in \texttt{code/}, with instructions in \texttt{README.md}.
\item \textbf{Data:} MPEG-7 2D images, two synthetic 3D surfaces, and the
Stanford Bunny reconstructed real scan geometry are included. Data provenance and Stanford Bunny dataset credits appear in \path{data/3d/README-data.txt}.
\item \textbf{Environment:} Python dependencies listed in \texttt{requirements.txt}; CPU execution is enough.
\item \textbf{Single-command main run:} From the project root, run \texttt{bash run.sh}. On Windows PowerShell use \texttt{python run\_all.py}. Both run the 2D experiment, procedural 3D and Bunny full-overlap experiments,
plus both 75\% source-coverage partial-overlap experiments and table generation.
\item \textbf{Extra experiments:} The robustness, parameter sweep, validation, and starting-pose diagnostics have separate scripts documented in the README.
\item \textbf{Single-case benchmark (optional):} Run \texttt{python benchmark\_single\_case.py} separately for a one-case timing and accuracy demonstration. Results are saved under \path{results/single_case_benchmark/}. This is not a repeated speed test.
\item \textbf{Random seeds:} 2D main seed 42. 3D data seed 2026; 3D experiment seed 20261112, with fixed shape, outlier, and trial offsets.
\item \textbf{Expected outputs:} 2D CSVs and figures; 3D results under
\path{results/3d/}; Bunny results under \path{results/3d_bunny/}.
Tables~\ref{tab:allrmse}, \ref{tab:3d}, and \ref{tab:bunny} give expected values.
\item \textbf{Pseudocode:} Separate \texttt{algorithm.txt} file in the project folder.
\item \textbf{Hardware and runtime:} Fill in the measured CPU, RAM, and running time above before submission.
\end{itemize}

\begin{thebibliography}{9}
\bibitem{liu2014}
J. Liu, S. Du, C. Zhang, J. Zhu, K. Li, and J. Xue,
``Probability Iterative Closest Point Algorithm for Position Estimation,''
\emph{2014 IEEE 17th International Conference on Intelligent Transportation Systems},
pp. 458--463, 2014.
\url{https://doi.org/10.1109/ITSC.2014.6957732}.

\bibitem{stanfordrepo}
Stanford University Computer Graphics Laboratory,
\emph{Stanford 3D Scanning Repository}; source of the scanned and reconstructed Bunny geometry used in our 3D tests.
\url{https://graphics.stanford.edu/data/3Dscanrep/}.

\bibitem{dataset}
MPEG-7 Core Experiment CE-Shape-1 Test Set, Part B,
\emph{Dataset information and download references},
Universidad del Pa\'is Vasco / Euskal Herriko Unibertsitatea.
\url{https://www.ehu.eus/ccwintco/index.php/MPEG-7_Core_Experiment_CE-Shape-1_Test_Set._Benchmarking_image_database_for_shape_recognition_techniques}.
\end{thebibliography}

\end{document}
